\hypertarget{appendices}{%
\section{9. Appendices}\label{appendices}}

\begin{quote}
Optional per the template; included here because the volume of
configuration data and the centrality of the handbook make them useful
to the reader.
\end{quote}

The order of appendices follows their importance to the reader. Appendix
A documents the handbook itself --- the durable, externally-hosted
artefact that this thesis writes \emph{about} --- and is intentionally
placed first. Appendix B is a glossary of acronyms used throughout the
thesis. Appendices C through G hold the configuration data, inventories,
diagrams, and governance material that supports the main text. Where the
same artefact is also documented in the companion thesis {[}Motje,
2026{]}, a cross-reference is given instead of a verbatim copy.

\begin{center}\rule{0.5\linewidth}{0.5pt}\end{center}

\hypertarget{appendix-a-the-online-handbook}{%
\subsection{Appendix A --- The online
Handbook}\label{appendix-a-the-online-handbook}}

The most important deliverable of this thesis is not contained in this
PDF. It is the \textbf{AUCOOP Community Network Handbook}, an online,
versioned, openly-licensed handbook hosted on GitHub and continuously
maintained by the AUCOOP association.

\hypertarget{a.1-where-to-find-it}{%
\subsubsection{A.1 Where to find it}\label{a.1-where-to-find-it}}

\begin{itemize}
\tightlist
\item
  \textbf{Web (canonical, evolving):} the latest version of the handbook
  is published as a static website built from the \texttt{main} branch
  of the source repository. The site is the recommended entry point for
  readers who want to browse, search, or follow internal cross-links.
\item
  \textbf{PDF (offline, expedition copy):} a self-contained PDF is
  published as a release artefact alongside each tagged version of the
  handbook, intended for download before travelling to a site with
  limited or unreliable connectivity (see §3.C.4).
\item
  \textbf{Source repository:}
  \url{https://github.com/aucoop/Community-Network-Handbook}
\item
  \textbf{Thesis-pinned snapshot:} the handbook content described in §3
  corresponds to branch \texttt{dev\_mj\_thesis} at commit
  \texttt{a5fc80b26368fe6cf1475994a12cc8f1967b84f0}, which collects the
  four feature branches merged for the purpose of this thesis (see
  Appendix E).
\end{itemize}

\hypertarget{a.2-structure-recap}{%
\subsubsection{A.2 Structure recap}\label{a.2-structure-recap}}

The handbook is organised in four top-level chapters with a strict
1-to-1 mapping between Chapter 2 (the imaginary use-case story) and
Chapter 3 (the recipe guide):

\begin{itemize}
\tightlist
\item
  \textbf{Chapter 1 --- Introduction.} Editorial purpose, audience, and
  how to read the handbook.
\item
  \textbf{Chapter 2 --- Imaginary Use Case.} A fictional community
  network told as a story, with one section per challenge, in narrative
  second-person voice with italic question titles.
\item
  \textbf{Chapter 3 --- Guide.} Step-by-step recipes, one folder per
  technology (\texttt{Network-Planning}, \texttt{IP-Addressing},
  \texttt{Wireless-Mesh}, \texttt{Antennas}, \texttt{Power-and-UPS},
  \texttt{Laptop-Deployment}, \ldots).
\item
  \textbf{Chapter 4 --- Real Use Cases.} Concrete deployments used as
  case studies (Namibia 2026 is the first such entry).
\end{itemize}

The structural rules and information architecture are described in
detail in §3.C.2. Rule files governing voice, structure, and navigation
conventions are reproduced in Appendix G.

\hypertarget{a.3-contributing}{%
\subsubsection{A.3 Contributing}\label{a.3-contributing}}

The handbook accepts contributions from anyone via the standard GitHub
pull-request workflow. New contributors are expected to read the
contributing guide and the four \texttt{.opencode/rules/*.md} files (see
Appendix G) before editing. The lightweight governance contract is
described in §3.C.6.

\hypertarget{a.4-how-to-cite-the-handbook}{%
\subsubsection{A.4 How to cite the
handbook}\label{a.4-how-to-cite-the-handbook}}

When citing the handbook, please cite the \textbf{specific commit or
release tag} consulted, not the moving \texttt{main} branch, so that the
reference remains reproducible:

\begin{quote}
AUCOOP. \emph{Community Network Handbook}. Version
\texttt{dev\_mj\_thesis\ @\ a5fc80b}.
\url{https://github.com/aucoop/Community-Network-Handbook}. Accessed
YYYY-MM-DD.
\end{quote}

The handbook should also be cited in conjunction with this thesis and
its companion {[}Motje, 2026{]}, since the three artefacts together
constitute the full record of the project.

\begin{center}\rule{0.5\linewidth}{0.5pt}\end{center}

\hypertarget{appendix-b-glossary-of-acronyms}{%
\subsection{Appendix B --- Glossary of
acronyms}\label{appendix-b-glossary-of-acronyms}}

The glossary lists acronyms and abbreviations used in the main text. In
the LaTeX build it is generated by the \texttt{acro} package; this
Markdown appendix is the editorial source.

\begin{longtable}[]{@{}
  >{\raggedright\arraybackslash}p{(\columnwidth - 2\tabcolsep) * \real{0.5000}}
  >{\raggedright\arraybackslash}p{(\columnwidth - 2\tabcolsep) * \real{0.5000}}@{}}
\toprule\noalign{}
\begin{minipage}[b]{\linewidth}\raggedright
Acronym
\end{minipage} & \begin{minipage}[b]{\linewidth}\raggedright
Expansion
\end{minipage} \\
\midrule\noalign{}
\endhead
\bottomrule\noalign{}
\endlastfoot
ADSL & Asymmetric Digital Subscriber Line \\
AP & Access Point \\
AUCOOP & Associació Universitària per la Cooperació (UPC student
association) \\
BIOS & Basic Input/Output System \\
BOM & Bill of Materials \\
CCD & Centre de Cooperació per al Desenvolupament (UPC development
cooperation centre) \\
CSV & Comma-Separated Values \\
dBm & decibel-milliwatt \\
DFS & Dynamic Frequency Selection \\
DHCP & Dynamic Host Configuration Protocol \\
DNS & Domain Name System \\
DRBL & Diskless Remote Boot in Linux \\
EFI & Extensible Firmware Interface \\
ESP & EFI System Partition \\
ESSID & Extended Service Set Identifier \\
ext4 & Fourth Extended Filesystem \\
GPT & GUID Partition Table \\
GRUB & GRand Unified Bootloader \\
HDD & Hard Disk Drive \\
ICT & Information and Communication Technology \\
IP & Internet Protocol \\
IPAM & IP Address Management \\
ISP & Internet Service Provider \\
LAN & Local Area Network \\
LEO & Low Earth Orbit \\
LTS & Long-Term Support \\
LuCI & OpenWrt web user interface \\
MAC & Media Access Control \\
MET & Master in Telecommunications Engineering \\
NFS & Network File System \\
NGO & Non-Governmental Organisation \\
NUT & Network UPS Tools \\
NVMe & Non-Volatile Memory Express \\
OS & Operating System \\
PoE & Power over Ethernet \\
PtP & Point-to-Point \\
PXE & Preboot eXecution Environment \\
RFC & Request for Comments \\
SAE & Simultaneous Authentication of Equals (WPA3) \\
SSD & Solid-State Drive \\
TFTP & Trivial File Transfer Protocol \\
UEFI & Unified Extensible Firmware Interface \\
UPC & Universitat Politècnica de Catalunya \\
UPS & Uninterruptible Power Supply \\
VPN & Virtual Private Network \\
WAN & Wide Area Network \\
WEEE & Waste Electrical and Electronic Equipment \\
WPA2 / WPA3 & Wi-Fi Protected Access (versions 2 and 3) \\
\end{longtable}

\begin{center}\rule{0.5\linewidth}{0.5pt}\end{center}

\hypertarget{appendix-c-bill-of-materials}{%
\subsection{Appendix C --- Bill of
materials}\label{appendix-c-bill-of-materials}}

Detailed BOM for the deployments described in §4: manufacturer part
numbers, suppliers, prices at time of purchase, and equivalent
substitutions noted where the part is not locally available.

To be filled at §4 / §5 writing time from the project purchasing record.

\hypertarget{appendix-d-pxe-server-configuration-files}{%
\subsection{Appendix D --- PXE server configuration
files}\label{appendix-d-pxe-server-configuration-files}}

Verbatim copies of the configuration files referenced in §3.B:

\begin{itemize}
\tightlist
\item
  \texttt{/etc/dhcp/dhcpd.conf}
\item
  \texttt{/etc/default/isc-dhcp-server}
\item
  \texttt{/etc/default/tftpd-hpa}
\item
  \texttt{/etc/exports}
\item
  \texttt{/tftpboot/nbi\_img/grub/grub.cfg}
\item
  \texttt{/home/partimag/auto-restore.sh}
\end{itemize}

Source of truth: \texttt{3-Guide/Laptop-Deployment/index.md} in the
handbook repository.

\hypertarget{appendix-e-repository-branch-and-commit-map}{%
\subsection{Appendix E --- Repository branch and commit
map}\label{appendix-e-repository-branch-and-commit-map}}

Snapshot of the handbook's git state at thesis submission time:

\begin{itemize}
\tightlist
\item
  \textbf{Pinned branch:} \texttt{dev\_mj\_thesis}
\item
  \textbf{Pinned commit:}
  \texttt{a5fc80b26368fe6cf1475994a12cc8f1967b84f0}
\item
  \textbf{Base branch:} \texttt{develop}
\item
  \textbf{Merged feature branches} (in merge order):

  \begin{itemize}
  \tightlist
  \item
    \texttt{docs/laptop-deployment}
  \item
    \texttt{WIP-docs/ip}
  \item
    \texttt{docs/planning}
  \item
    \texttt{fix/open-points}
  \end{itemize}
\item
  \textbf{Per-chapter authorship table} extracted from
  \texttt{git\ shortlog} at the pinned commit (to be inserted at
  submission time).
\end{itemize}

The Namibia case study in \texttt{docs/4-Real-Use-Cases/4.1-Namibia/}
was unified from the wireless-network and laptop-deployment narratives
during the merge of \texttt{fix/open-points}.

\hypertarget{appendix-f-network-and-deployment-diagrams-namibia}{%
\subsection{Appendix F --- Network and deployment diagrams
(Namibia)}\label{appendix-f-network-and-deployment-diagrams-namibia}}

\begin{itemize}
\tightlist
\item
  Pre-deployment topology of N Mutschuana Primary School
\item
  Post-deployment topology with mesh routers and IP plan annotations
\item
  IP addressing plan visualisation
\item
  PXE deployment subnet diagram
\end{itemize}

Source: \texttt{4-Real-Use-Cases/4.1-Namibia/images/} and
\texttt{3-Guide/Laptop-Deployment/images/} in the handbook repository.

\hypertarget{appendix-g-handbook-governance-materials}{%
\subsection{Appendix G --- Handbook governance
materials}\label{appendix-g-handbook-governance-materials}}

Reproduced verbatim from the pinned commit of the handbook repository so
the reader of this thesis can audit the conventions described in §3.C
without leaving the document:

\begin{itemize}
\tightlist
\item
  \texttt{AGENTS.md} --- repository-level OpenCode rules and build
  commands.
\item
  \texttt{.opencode/rules/general.md} --- general handbook conventions.
\item
  \texttt{.opencode/rules/chapter2-story.md} --- narrative voice for
  Chapter 2.
\item
  \texttt{.opencode/rules/chapter3-guide.md} --- recipe structure for
  Chapter 3.
\item
  \texttt{.opencode/rules/mkdocs-nav.md} --- navigation-manifest
  invariants.
\item
  \texttt{.opencode/agents/\{writer,reviewer,structure,diagrams,consistency\}.md}
  --- subagent definitions.
\item
  \texttt{.opencode/commands/\{new-section,review-chapter,add-diagram,audit,guide-from-steps\}.md}
  --- custom command definitions.
\end{itemize}

These materials are the AUCOOP-specific implementation of the governance
contract of §3.C.6.
